\documentclass[11pt]{article}

\usepackage[utf8]{inputenc}
\usepackage[T1]{fontenc}
\usepackage{amsmath,amssymb}
\usepackage{booktabs}
\usepackage{graphicx}
\graphicspath{{./}}
\usepackage{hyperref}
\usepackage[margin=1in]{geometry}
\usepackage{url}

\hypersetup{colorlinks=true, linkcolor=blue, citecolor=blue, urlcolor=blue}

\title{Ordering Sensitivity in Subgraph-Isomorphism Layout Search:\\
A Characterization of Catastrophic Failure Regions in\\
Quantum Circuit Transpilation}

\author{TN Holdings\\
\texttt{love.os.architect@proton.me}\\
\texttt{https://github.com/TN-Holdings-LLC/psf-zero}}

\date{\today}

\begin{document}
\maketitle

\begin{abstract}
Qubit layout selection in quantum circuit transpilation is commonly
formulated as subgraph isomorphism and solved with VF2-family search.
We report that Qiskit's \texttt{VF2Layout} pass exhibits a sharp,
reproducible failure region: on interaction graphs that exactly
saturate a device's maximum matching, median search time rises from
$34.1$\,ms to $9{,}236$\,ms---a factor of $271$---and terminates by
declaring no solution exists, on instances where a valid embedding
demonstrably does exist. The failure is not a property of
subgraph-isomorphism search as such. Reading Qiskit's current Rust
source, we find \texttt{VF2Layout} does not call the public
\texttt{rustworkx} mapper but a separate implementation with a
hardcoded VF2++ node-ordering strategy, applied unconditionally. On
the same 26 saturated instances, a single call to
\texttt{rustworkx.vf2\_mapping} with a fixed, unshuffled node order
(\texttt{id\_order=True}) finds a valid embedding in all 26 cases, at a
median of $0.0292$\,ms---roughly five orders of magnitude faster than
the failing search on the 19 instances where Qiskit reports no
solution. We show by faithful simulation of Qiskit's own
\texttt{Vf2ppSorter} that VF2++ ordering defers the highly symmetric,
degree-one region of these interaction graphs to the end of its
traversal, and we report that this mechanism alone does not account for
the observed instance-level divergence: graph-isomorphic instance pairs
differing only in which physical qubits host which component produce
opposite outcomes, reproducibly. A permutation-based analysis across
all 325 instance pairs finds that no graph-level feature we computed
explains more than $4.5\%$ of outcome variance, with none surviving
correction for multiple comparisons. We report a cross-SDK comparison
in which Cirq's placement search exceeds a 10-second budget across a
substantially wider occupancy range than Qiskit's failure region, and
topology results on brick and heavy-hex lattices. All predictions were
pre-registered before measurement; all raw data, pre-registrations, and
analysis scripts are public.
\end{abstract}

\section{Introduction}

Mapping a logical quantum circuit onto physical hardware requires
choosing an assignment of logical qubits to physical qubits such that
every two-qubit interaction lands on a physically coupled pair. When a
perfect assignment exists---that is, when the circuit's interaction
graph embeds as a subgraph of the device's coupling graph---finding it
is a subgraph isomorphism problem, and both Qiskit and other major
transpilation frameworks attempt exactly this before falling back to
heuristic routing with SWAP insertion.

Subgraph isomorphism is NP-complete in general, and practical solvers
accordingly impose call and time budgets. What is less appreciated is
how sharply the practical difficulty of these instances depends on
factors that are invisible to the problem statement: not merely the
size of the graphs, but the specific node-traversal order the solver
happens to use, and---as we show---which physical qubits happen to host
which part of an otherwise identical interaction graph.

This paper characterizes a specific, reproducible failure region in
Qiskit's \texttt{VF2Layout} pass, identifies the ordering strategy
responsible, and demonstrates that an extremely simple alternative
resolves every instance in the region. We are deliberate about the
boundary between what we measured and what we can explain: a
substantial part of this paper reports mechanisms we ruled out and
regularities we failed to find, because those negative results
constrain the problem more sharply than the positive ones alone.

\subsection{Contributions}

\begin{enumerate}
\item \textbf{Characterization.} A reproducible failure region in
Qiskit's \texttt{VF2Layout}: on saturated instances, median search time
of $9{,}236$\,ms (range $8{,}815$--$11{,}442$\,ms) terminating in a
``no solution exists'' verdict, against $34.1$\,ms for non-failing
instances in the same family---on instances where a valid embedding
provably exists (Section~\ref{sec:cliff}).

\item \textbf{Mechanism, partially.} Qiskit's \texttt{VF2Layout} uses a
custom implementation with hardcoded VF2++ ordering, verified from
source, not the public \texttt{rustworkx} mapper. A faithful
reimplementation of its \texttt{Vf2ppSorter} shows the ordering defers
the symmetric degree-one region of these instances to the end of the
traversal (Section~\ref{sec:mechanism}).

\item \textbf{A minimal alternative.} A single
\texttt{rustworkx.vf2\_mapping} call with \texttt{id\_order=True}
resolves all 26 instances at a median of $0.0292$\,ms
(Section~\ref{sec:ordering}).

\item \textbf{Negative results that constrain the problem.}
Graph-isomorphic instance pairs diverge reproducibly; no computed
graph-level feature explains more than $4.5\%$ of outcome variance;
the ordering mechanism alone does not account for instance-level
divergence (Section~\ref{sec:unexplained}).

\item \textbf{Methodology.} Every prediction in this work was
pre-registered in a timestamped document before the corresponding
measurement, including predictions that were subsequently falsified.
Several of our own earlier conclusions were corrected in the course of
the investigation and the corrections are part of the public record
(Section~\ref{sec:method}).
\end{enumerate}

\section{Background}

\subsection{Layout as subgraph isomorphism}

Let $G_I = (V_I, E_I)$ be a circuit's interaction graph, with a vertex
per logical qubit and an edge per interacting pair, and let
$G_P = (V_P, E_P)$ be the device coupling graph. A \emph{perfect
layout} is an injection $\phi : V_I \to V_P$ with
$(\phi(u), \phi(v)) \in E_P$ for every $(u,v) \in E_I$. When such a
$\phi$ exists, no SWAP insertion is required.

We parameterize instances by \emph{spare qubits}: the number of
physical qubits not used by the interaction graph. At
$\mathrm{spare}=0$ the interaction graph uses every physical qubit,
leaving the search no freedom to place a component anywhere unoccupied.

\subsection{VF2 and VF2++}

VF2~\cite{vf2} explores a search tree of partial mappings, extending
one vertex at a time and pruning by local feasibility. Its practical
performance depends heavily on the order in which pattern vertices are
considered. VF2++~\cite{vf2pp} improves on this with a
structure-informed ordering: roughly, process vertices in an order that
maximizes connectivity to the already-mapped set, breaking ties by
degree. On most instance families this ordering is a substantial
improvement, which is why it is widely adopted.

\section{Methodology}
\label{sec:method}

\subsection{Software environment}

All measurements reported in this paper were taken on
\texttt{AMD64 Family 25 Model 80}, with Qiskit \texttt{2.5.2} and
\texttt{rustworkx} \texttt{0.18.1}. Section~\ref{sec:cliff}'s primary
measurement (Table~\ref{tab:cliff}) used Python \texttt{3.10.11};
Section~\ref{sec:ordering}'s ordering-sensitivity measurements
(Table~\ref{tab:topology}, and the $0.0292$\,ms figure) used Python
\texttt{3.11.9}. We report both explicitly rather than a single value,
since the two measurements were taken in separate sessions and we did
not verify the Python patch version was held constant between them; we
have no evidence this affects any reported result, and note it for
completeness rather than as a caveat we can currently resolve.

\subsection{Pre-registration}

Every experiment reported here was preceded by a pre-registration
document stating the design and the predicted outcomes, including
explicit statements of what result would falsify the hypothesis, written
and committed before the corresponding measurement was performed. These
documents are public and are cited individually below by their addendum
number.

This discipline was not cosmetic. Over the course of this
investigation, several attractive hypotheses were pre-registered and
then falsified by their own designated tests---a component-count
divisibility rule, a size-ratio threshold, and a locality-based
predictive model among them. Several of our own published intermediate
conclusions were subsequently corrected; in each case the correction is
recorded alongside the original rather than replacing it. We regard
this record as part of the contribution, since the central negative
result of Section~\ref{sec:unexplained} is only credible to the extent
that the positive hypotheses preceding it were genuinely tested rather
than abandoned quietly.

\subsection{Instance families}

Our primary family, \texttt{dominant\_size\_sweep}, is parameterized by
a single integer $D$. On $n=64$ qubits it comprises 17 disjoint
two-qubit edges (34 qubits) plus one path of $D$ qubits and one path of
$30-D$ qubits, for $D \in \{2, \ldots, 28\}$. Every member has zero
idle qubits and an interaction-graph maximum matching of exactly 32,
equal to that of the $8\times8$ grid---genuine saturation, verified
numerically per instance rather than assumed.

The family has a property we exploit throughout: for any $D$, the
instances at $D$ and $30-D$ are graph-isomorphic, differing only in
which physical qubit indices host which path. This gives us matched
pairs that any purely structural account must treat identically.

\subsection{Feasibility}

For every instance reported as a failure, we verify independently
(via \texttt{networkx} maximum matching on the interaction graph
against the coupling graph) that a valid embedding exists. Every
failure discussed in this paper is a failure to find an embedding that
does exist, not a correct report of infeasibility.

\section{The failure region}
\label{sec:cliff}

\subsection{Primary measurement}

Figure~\ref{fig:cliff} shows the qualitative shape of this failure
region on an independent, earlier measurement using this project's
original disjoint-pair circuit family (not
\texttt{dominant\_size\_sweep}) on the same $6\times7$ grid: VF2Layout
time drops by roughly two orders of magnitude the moment a single spare
qubit is introduced, and stays flat from spare$=1$ onward. We include
it for its visual clarity; the quantitative results in the rest of this
section use \texttt{dominant\_size\_sweep} exclusively, and the two
families are not compared to each other numerically anywhere in this
paper.

\begin{figure}[h]
\centering
\includegraphics[width=0.75\textwidth]{occupancy_cliff.pdf}
\caption{VF2Layout time versus spare qubits, $6\times7$ grid,
disjoint-pair circuits, optimization level 3, median of 18 runs per
point (3 seeds $\times$ 6 repeats), shaded region spans min--max.
Independent measurement from Table~\ref{tab:cliff}'s
\texttt{dominant\_size\_sweep} results; shown for illustration of the
failure region's shape, not as a quantitative source for this paper's
claims.}
\label{fig:cliff}
\end{figure}

Table~\ref{tab:cliff} summarizes \texttt{VF2Layout} at
\texttt{optimization\_level=3} on the $8\times8$ grid across the
\texttt{dominant\_size\_sweep} family, 3 seeds $\times$ 2 repeats per
instance, replicated in two independent batches.

\begin{table}[h]
\centering
\begin{tabular}{lrrr}
\toprule
Outcome & Instances & Median time & Range \\
\midrule
\texttt{nonexistent solution} & 13 & $9{,}236$\,ms & $8{,}815$--$11{,}442$\,ms \\
\texttt{solution found} & 5 & $34.1$\,ms & --- \\
\midrule
Ratio & & $271\times$ & \\
\bottomrule
\end{tabular}
\caption{Qiskit \texttt{VF2Layout}, $8\times8$ grid,
\texttt{dominant\_size\_sweep}, 18 instances with both batches
complete. Every ``nonexistent solution'' instance admits a valid
embedding.}
\label{tab:cliff}
\end{table}

Two features are worth emphasizing. First, the failing instances do not
merely take longer---they terminate with a verdict of no solution, on
instances where one exists. Downstream, this means the transpiler
proceeds to heuristic routing and inserts SWAP gates that were not
necessary. Second, the outcome is bimodal: instances land in one regime
or the other with no intermediate values, and the classification is
stable across independently replicated batches (18/18 instances agreed
between batches).

\subsection{Necessary conditions}

Two conditions are jointly necessary for the failure region in our
measurements. Introducing a single idle qubit ($\mathrm{spare}=1$)
eliminates it. Connecting the interaction graph's components into a
single connected structure also eliminates it. Neither alone is
sufficient: connected interaction graphs at $\mathrm{spare}=0$ do not
fail, and disconnected graphs with spare capacity do not fail.

\subsection{Cross-SDK comparison}

We ran an analogous occupancy sweep against Cirq's own placement search
on a $6\times7$ grid with a 10-second budget per instance. Results are
in Table~\ref{tab:cirq}.

\begin{table}[h]
\centering
\begin{tabular}{lrr}
\toprule
Spare qubits & Median elapsed & Timed out \\
\midrule
0 & $\geq 10.03$\,s & yes \\
2 & $\geq 10.03$\,s & yes \\
4 & $\geq 10.03$\,s & yes \\
6 & $\geq 10.03$\,s & yes \\
8 & $\geq 10.02$\,s & yes \\
16 & $2.31$\,s & no \\
24 & $1.40$\,s & no \\
\bottomrule
\end{tabular}
\caption{Cirq placement search, $6\times7$ grid. Values at
$\mathrm{spare} \leq 8$ are right-censored at the 10-second budget and
are lower bounds, not measurements. A valid embedding exists at every
row.}
\label{tab:cirq}
\end{table}

We stress the censoring: we cannot report how long Cirq would take,
only that it exceeds 10 seconds. The comparison we can make is
qualitative but notable in its own right---Cirq's difficult region
extends across $\mathrm{spare} = 0$ through $8$, substantially wider
than the $\mathrm{spare}=0$-only region we observe in Qiskit. Whatever
the underlying cause, it is not unique to one implementation, and its
extent differs between implementations.

\section{What Qiskit actually runs}
\label{sec:mechanism}

\subsection{Source-level verification}

A natural hypothesis is that Qiskit's \texttt{VF2Layout} delegates to
\texttt{rustworkx}'s public \texttt{vf2\_mapping}. Reading Qiskit's
current Rust source establishes that it does not. The relevant pass
imports only graph data structures from \texttt{rustworkx\_core}; the
search itself is \texttt{qiskit\_circuit::vf2}, a separate
implementation. Both the ``average'' and ``exact'' scoring variants
construct their search identically with respect to ordering:

\begin{verbatim}
let vf2 = vf2::Vf2::new(&interactions.graph, &coupling_graph,
                        vf2::Problem::Subgraph)
    .with_scoring(score, score)
    .with_restriction(vf2::Restriction::Decreasing(best_score))
    .with_vf2pp_ordering();
\end{verbatim}

\noindent
\texttt{with\_vf2pp\_ordering()} is called unconditionally in both. We
found no code path in this file that runs the search with an
alternative ordering, and no configuration option exposing one. This
distinction matters for the rest of the paper: the alternative we
report in Section~\ref{sec:ordering} is not a setting Qiskit could
toggle, but a different implementation.

\subsection{Simulating the ordering}

We reimplemented \texttt{Vf2ppSorter}'s ordering rule---select by
(connectivity to the already-ordered set, total degree, lowest index),
building a BFS tree per root---and ran it on our instances.

The result is unambiguous. In these interaction graphs, path-interior
vertices have degree 2 and bare-edge vertices have degree 1. The
ordering rule therefore places every path vertex before any bare-edge
vertex: for $D=10$ on $n=64$, the first bare-edge vertex appears at
position 30 of 64, exactly when the 30-qubit path allocation is
exhausted. The symmetric, mutually interchangeable region of the graph
is processed last, as a block, after all the structurally
distinguishable vertices are placed.

This is consistent with Qiskit's own release notes for version 2.2,
which reduced call and trial limits for this pass citing ``excessive
runtimes for highly symmetric trial circuits being mapped to large
coupling maps.'' The instances in our failure region are exactly that:
17 mutually interchangeable two-qubit components.

\subsection{Why this is only half a mechanism}

The ordering account explains why these instances are hard. It does not
explain which of them fail. For the isomorphic pair $D=10$ and $D=20$,
the total path allocation is 30 qubits either way, so the ordering
produces structurally parallel traversals: both process 30 path
vertices, then reach the bare-edge region at position 30. Yet $D=10$
succeeds and $D=20$ fails, reproducibly across independent batches.

Whatever distinguishes them must lie in the backtracking search's
behavior once it reaches the tied region, not in the initial ordering.
We did not simulate the full backtracking search, and we do not claim
to know the answer.

\section{Ordering sensitivity}
\label{sec:ordering}

\subsection{Result}

We ran a single \texttt{rustworkx.vf2\_mapping} call per instance with
\texttt{subgraph=True}, \texttt{induced=False}, and
\texttt{id\_order=True}---a fixed traversal in natural vertex order, with
no structure-informed heuristic at all---on the physical graph in its
unmodified numbering.

It finds a valid embedding in \textbf{26 of 26} instances, including
all 19 where Qiskit reports no solution. Median time: $0.0292$\,ms;
maximum across all instances: $0.108$\,ms.

Against the failing instances' median of $9{,}236$\,ms, this is a
factor of approximately $3.2 \times 10^{5}$. We report the comparison
with two caveats. It is a comparison between two different
implementations, not one implementation under two settings
(Section~\ref{sec:mechanism}); and the Qiskit figure is the time to
reach a wrong answer, so the speedup and the correctness improvement
are not separable.

\subsection{Topology generalization}

To test whether fixed ordering is merely lucky on square grids, we
evaluated it on two structurally different topologies. Because the
relevant comparison here is between orderings within a single
implementation, both arms use \texttt{rustworkx}: fixed order versus
the family of BFS- and degree-based orderings that a prior prototype of
ours had used, on the same instances.

\begin{table}[h]
\centering
\begin{tabular}{lrrrrr}
\toprule
Topology & $n$ & Fixed order & Time & Other orderings & Time \\
\midrule
brick $6\times7$ & 42 & found & $1.25$\,ms & none found & $875$\,ms \\
brick $8\times8$ & 64 & found & $1.08$\,ms & none found & $996$\,ms \\
heavy-hex $d{=}3$ & 19 & found & $0.25$\,ms & --- & --- \\
heavy-hex $d{=}5$ & 57 & found & $0.69$\,ms & --- & --- \\
\bottomrule
\end{tabular}
\caption{Fixed ordering versus structure-informed orderings, disjoint-pair
instances at saturation. On brick, the alternatives exhaust nine
orderings without success. Heavy-hex instances are saturated relative
to the lattice's own maximum matching, since its bipartite imbalance
makes qubit-count saturation unattainable.}
\label{tab:topology}
\end{table}

On brick lattices the contrast is stark: the fixed order succeeds in
about a millisecond on instances where nine structure-informed
orderings, given a substantially larger budget, find nothing at all.

\subsection{Interpretation}

We do not claim fixed ordering is superior in general; a
structure-informed heuristic that is better on typical instances and
worse on symmetric ones is an entirely reasonable engineering
trade-off, and VF2++ is widely adopted for good reason. What our
results establish is narrower and, we think, more useful: on saturated,
highly symmetric layout instances---a regime that quantum transpilation
specifically cares about, since maximizing device utilization is the
point---the sophisticated ordering is not merely slower but
qualitatively worse, and the cheapest possible alternative is not
merely competitive but dominant by orders of magnitude.

A practical implication follows directly. A transpiler could attempt a
budgeted fixed-order pass before its primary search at negligible cost
($\leq 0.11$\,ms in all our measurements) and fall through on failure,
gaining the entire failure region at no meaningful expense.

\section{What we could not explain}
\label{sec:unexplained}

This section reports negative results. We consider them a contribution
rather than a gap: they bound how much of this phenomenon is currently
understood, and they close off several natural explanations that a
reader would otherwise reasonably propose.

\subsection{Isomorphic instances diverge}

Of the 13 graph-isomorphic pairs $(D, 30-D)$ in our family, 5 diverge:
one member succeeds, the other fails, reproducibly. The remaining pairs
agree (7 fail on both sides; one is a degenerate construction).

We confirmed isomorphism directly. The pairs differ only in which
physical qubit indices host which path. This means \textbf{no account
based purely on the abstract structure of the interaction graph can be
complete}---physical placement, not graph structure alone, determines
the outcome.

One regularity survived every test we applied to it: in all 5 diverging
pairs, the member with the \emph{smaller} dominant path succeeds. We
report this without a proposed mechanism, and note that our record
contains several patterns of comparable apparent strength that
subsequent tests falsified.

\subsection{A ratio threshold, falsified}

Having observed divergence at size ratios $1.5$ and $2.0$ and agreement
at $2.75$ and $5.0$, we pre-registered a threshold hypothesis and
tested all 13 pairs. The result was non-monotonic: ratios $6.5$ and
$3.29$ diverge while $4.0$ and $1.73$ agree. The hypothesis is
falsified.

\subsection{No predictive feature}

Rather than continue testing candidate variables one at a time---a
process that can fail indefinitely without concluding anything---we
reframed the question as a measurable one: how much do the features we
can compute actually explain?

We computed pairwise distances across all $\binom{26}{2} = 325$
instance pairs under three metrics and correlated each against outcome
agreement, with significance assessed by permutation test
(5,000 shuffles).

\begin{table}[h]
\centering
\begin{tabular}{lrrrl}
\toprule
Distance metric & $r$ & $p$ & $r^2$ & Survives Bonferroni \\
\midrule
Dominant size & $-0.151$ & $0.032$ & $2.3\%$ & no \\
Size ratio & $-0.212$ & $0.049$ & $4.5\%$ & no \\
Max component size & $-0.017$ & $0.354$ & $0.0\%$ & no \\
\bottomrule
\end{tabular}
\caption{Predictive power of computed graph features. Correction for
three tests requires $p < 0.0167$.}
\label{tab:mantel}
\end{table}

No metric survives correction, and the largest effect explains under
$5\%$ of variance. We state the conclusion precisely: \emph{the
graph-level features we computed have little predictive power over this
outcome.} This is not a claim that no predictive feature exists. Our
instance family has essentially one independent numeric parameter, so
no correlational analysis of it can detect structure depending on
information that parameter does not carry---including, notably, the
physical placement that Section 6.1 shows does matter.

\subsection{Summary of the boundary}

We can state when the failure region occurs (saturation, disconnected
interaction graph, sufficient symmetric structure), what ordering
Qiskit uses and how it treats these instances, and that a trivial
alternative resolves every instance we tested. We cannot predict which
specific instances fail, explain why isomorphic instances diverge, or
account for the ``smaller succeeds'' regularity.

\section{Related work}

VF2~\cite{vf2} and VF2++~\cite{vf2pp} are the algorithmic basis for the
layout passes examined here. Qiskit's transpiler
architecture~\cite{qiskit} places \texttt{VF2Layout} before heuristic
routing precisely so that perfect layouts, when they exist, are found
and SWAP insertion avoided; SABRE~\cite{sabre} is the heuristic
fallback when they are not.

Ordering sensitivity in backtracking search is well known in the
constraint-satisfaction literature, and symmetric instances are a
recognized worst case for subgraph isomorphism generally. Our
contribution is not the observation that ordering matters, but the
measurement of its magnitude in a specific, deployed, practically
important setting---and the finding that the direction of the effect
reverses relative to the general case, with the sophisticated heuristic
losing to the naive one by five orders of magnitude on the instances
that matter most for device utilization.

\section{Conclusion}

Layout search in quantum transpilation has a failure region that is
sharp, reproducible, and consequential: on saturated instances, Qiskit's
\texttt{VF2Layout} spends seconds to reach an incorrect verdict of
infeasibility, on instances that a fixed-order search resolves in tens
of microseconds. The mechanism is at least partly the node ordering:
VF2++'s structure-informed traversal defers the symmetric region of
these instances to the end, where the search then struggles.

We have been explicit about the limits of this account. It does not
predict which instances fail; isomorphic instances diverge on physical
placement alone; no feature we computed explains meaningful outcome
variance. We report these as results rather than omissions, because
they are what currently bounds understanding of the phenomenon, and
because the practical recommendation---attempt a fixed-order pass first,
at negligible cost---does not depend on resolving them.

\section*{Data and code availability}

All raw measurement data, pre-registration documents, analysis scripts,
and the complete chronological record including superseded conclusions
and their corrections are available at
\url{https://github.com/TN-Holdings-LLC/psf-zero}.

\begin{thebibliography}{9}

\bibitem{vf2}
L.~P. Cordella, P.~Foggia, C.~Sansone, and M.~Vento,
\emph{A (sub)graph isomorphism algorithm for matching large graphs},
IEEE Transactions on Pattern Analysis and Machine Intelligence,
26(10):1367--1372, 2004.

\bibitem{vf2pp}
A.~Juttner and P.~Madarasi,
\emph{VF2++---An improved subgraph isomorphism algorithm},
Discrete Applied Mathematics, 242:69--81, 2018.

\bibitem{qiskit}
Qiskit contributors,
\emph{Qiskit: An Open-source Framework for Quantum Computing},
2023. \url{https://qiskit.org}

\bibitem{sabre}
G.~Li, Y.~Ding, and Y.~Xie,
\emph{Tackling the qubit mapping problem for NISQ-era quantum devices},
ASPLOS '19, pp.~1001--1014, 2019.

\end{thebibliography}

\end{document}
